\documentclass[12pt]{article}
\usepackage[utf8]{inputenc}
\usepackage[T1]{fontenc}
\usepackage{hyperref}
\title{Waskia Language Resources for K-12 Curriculum Development}
\author{Compiled by Claude Sonnet 5.5 (claude-sonnet-5-5) for Andr\'as Kornai}
\date{}
\begin{document}
\maketitle

\section{Introduction}
Waskia (ISO 639-3 wsk; Glottocode wask1241) is a Trans-New Guinea language of the Madang branch, spoken in Madang Province, Papua New Guinea, with its centre on Karkar Island. Estimates are about 20,000 speakers; one Bible preface gives about 35,000. Joshua Project describes Waskia as in active daily use. Searches found no school curriculum, primer or textbook in Waskia. The material below is mainly reference and religious material.

\section{Primary Educational Resources}
No Waskia primers, textbooks or school programmes turned up in search results. A curriculum would have to be built from scratch, drawing on the reference and scripture materials below.

\section{Linguistic and Reference Materials}
\begin{itemize}
\item Wikipedia overview of the language: \url{https://en.wikipedia.org/wiki/Waskia_language}
\item Omniglot page on Waskia: \url{https://www.omniglot.com/writing/waskia.htm}
\item Ross and Paol, \emph{A Waskia grammar sketch and vocabulary} (1978, Australian National University), reported to be on the Internet Archive. No direct URL was seen in the search results.
\item Barker and Lee, \emph{A tentative phonemic statement of Waskia} (SIL). No direct URL was seen in the search results.
\end{itemize}

\section{Religious and Community Resources}
\begin{itemize}
\item The 1985 Waskia New Testament, with Genesis, Proverbs and portions of Psalms, at the Digital Bible Society: \url{https://dbs.org/bibles/WSKPNG}
\item \emph{Kaem ko den}, a Waskia text on eBible.org: \url{https://ebible.org/study/content/texts/wsk/FR0.html}
\item Joshua Project people-group profile: \url{https://joshuaproject.net/people_groups/15883}
\end{itemize}

\section{Digital and Audio Resources}
The Endangered Languages Documentation Programme collection ``Gildipasi language project: tumbuna stories and tumbuna knowledge'' covers Waskia, Barem and Manep in Madang Province: \url{https://eldp.access.preservica.com/dk0559/}

\section{Curriculum Development Context}
The language is still widely spoken, including by children in the Karkar Island setting described by Joshua Project, but no schooling in Waskia was found. The documented grammar, the Bible translation and the ELDP recordings could support a mother-tongue literacy programme. Any such programme would need community and local education authority involvement. Little else was found, and the gaps above should be filled by fieldwork or by contacting Papua New Guinea education authorities.

\end{document}